\documentclass[12pt]{article}
\usepackage{color}
\usepackage{amsmath}
\usepackage{arxiv}
\usepackage{lyu}
\usepackage{xcolor}
\usepackage{pgfplots}
\usetikzlibrary{calc}
\pgfplotsset{compat=1.18}
\pgfplotsset{colormap={violet}{rgb255=(25,25,122) rgb255=(238,140,238) color=(white)}}
\usepackage[utf8]{inputenc} % allow utf-8 input
\usepackage[T1]{fontenc}    % use 8-bit T1 fonts
\usepackage{hyperref}       % hyperlinks
\usepackage{url}            % simple URL typesetting
\usepackage{booktabs}       % professional-quality tables
\mathchardef\mhyphen="2D
\usepackage{amsfonts}       % blackboard math symbols
\usepackage{nicefrac}       % compact symbols for 1/2, etc.
\usepackage{amsthm}
\newtheorem{theorem}{Theorem}
\newtheorem{proposition}{Proposition}[section]
\newtheorem{corollary}{Corollary}[theorem]
\newtheorem{lemma}[theorem]{Lemma}
\usepackage{microtype}      % microtypography
\usepackage{lipsum}
\usepackage{comment}
\usepackage{subcaption}
\usepackage{graphicx}
\usepackage{soul}  
\title{proof}
\begin{document}
We first isolate the error introduced by the operator splitting itself, i.e.\ the error of replacing the exact flow of \eqref{equ:normalized_diffusion_reaction} by the Lie--Trotter recursion \eqref{eq:lie-trotter}, at the continuous
(PDE) level and independently of the particle and spectral discretizations.

Recall the operators $\mathcal A$ and $\mathcal B$ defined in
\eqref{eq:operator-A}--\eqref{eq:operator-B}. Let $\Phi^{\tau}$ denote
the solution flow of $\mathcal A+\mathcal B$ over a step $\tau$, and
$\Phi^{\tau}_{\mathcal A},\Phi^{\tau}_{\mathcal B}$ the flows of the
sub-equations $\partial_t\rho=\mathcal A(\rho)$ and
$\partial_t\rho=\mathcal B(\rho)$. One Lie--Trotter step is
$\mathcal S_{\tau}:=\Phi^{\tau}_{\mathcal B}\circ\Phi^{\tau}_{\mathcal A}$, and the
time-discrete split solution is $\rho^{0}=\rho_0$,
$\rho^{n+1}=\mathcal S_{\tau}\rho^{n}$, $t_n=n\tau$.
We denote by $\mathcal A'(\rho),\mathcal B'(\rho)$ the Fréchet derivatives of
$\mathcal A,\mathcal B$ at $\rho$, and define the nonlinear Lie bracket
\begin{equation}
[\mathcal A,\mathcal B](\rho)
:=\mathcal A'(\rho)\,\mathcal B(\rho)-\mathcal B'(\rho)\,\mathcal A(\rho).
\label{eq:bracket}
\end{equation}


\begin{assumption}\label{ass:regularity}
   $\tilde a\in C^{2}_{b}(\mathbb R;\mathbb R^{d})$,
$\tilde r\in C^{2}_{b}(\mathbb R)$, and the exact solution of
\eqref{eq:rho-pde} satisfies
$\rho\in C\big([0,T];H^{p}(\mathbb T^{d})\big)$ for some integer
$p>\tfrac d2+2$, with
$\sup_{0\le t\le T}\|\rho(t)\|_{H^{p}}\le R$.
\end{assumption}

\begin{lemma}[Local splitting error]\label{lem:local-split}
Under Assumption~\ref{ass:regularity}, there exist $\tau_0>0$ and a constant
$C_{\mathrm L}=C_{\mathrm L}\big(R,D,\|\tilde a\|_{C^{2}_b},\|\tilde r\|_{C^{2}_b},d\big)$
such that, for every $\phi$ with $\|\phi\|_{H^{p}}\le R$ and every
$\tau\in(0,\tau_0]$,
\[
\big\|\Phi^{\tau}\phi-\mathcal S_{\tau}\phi\big\|_{L^{2}}\le C_{\mathrm L}\,\tau^{2}.
\]
More precisely, the defect is governed to leading order by the bracket
\eqref{eq:bracket},
\begin{equation}
\Phi^{\tau}\phi-\mathcal S_{\tau}\phi
=\tfrac{\tau^{2}}{2}\,[\mathcal A,\mathcal B](\phi)+O(\tau^{3}).
\label{eq:leading}
\end{equation}
\end{lemma}

\begin{proof}
For a smooth flow $\Phi^{t}_{\mathcal F}$ generated by $\mathcal F$, Taylor expansion
in $t$ gives
$\Phi^{t}_{\mathcal F}\phi=\phi+t\,\mathcal F(\phi)
+\tfrac{t^{2}}{2}\mathcal F'(\phi)\mathcal F(\phi)+O(t^{3})$.
Applying this to $\Phi^{\tau}_{\mathcal A}$ and then $\Phi^{\tau}_{\mathcal B}$, and
expanding $\mathcal B$ and $\mathcal B'$ about $\phi$, one obtains
\[
\mathcal S_{\tau}\phi=\phi+\tau(\mathcal A+\mathcal B)\phi
+\tfrac{\tau^{2}}{2}\big(\mathcal A'\mathcal A+\mathcal B'\mathcal B\big)\phi
+\tau^{2}\,\mathcal B'\mathcal A\,\phi+O(\tau^{3}),
\]
while the exact flow of $\mathcal C:=\mathcal A+\mathcal B$ satisfies
\[
\Phi^{\tau}\phi=\phi+\tau\,\mathcal C\phi
+\tfrac{\tau^{2}}{2}\big(\mathcal A'\mathcal A+\mathcal A'\mathcal B
+\mathcal B'\mathcal A+\mathcal B'\mathcal B\big)\phi+O(\tau^{3}).
\]
Subtracting cancels the $O(\tau)$ and the symmetric $O(\tau^{2})$ terms and
leaves $\tfrac{\tau^{2}}{2}\big(\mathcal A'\mathcal B-\mathcal B'\mathcal A\big)\phi
+O(\tau^{3})$, which is \eqref{eq:leading}.The leading-order term in \eqref{eq:leading} satisfies
$\big\|\tfrac{\tau^{2}}{2}[\mathcal A,\mathcal B](\phi)\big\|_{L^{2}}
\le C_{\mathrm L}\tau^{2}$, since Assumption~\ref{ass:regularity} bounds
$\mathcal A(\phi),\mathcal B(\phi)$ and the bracket \eqref{eq:bracket} in $L^{2}$
uniformly for $\|\phi\|_{H^{p}}\le R$, via the embedding
$H^{p-2}\hookrightarrow L^{\infty}$ and the algebra property of $H^{p-2}$.
The third-order remainder in \eqref{eq:leading} admits the integral form of
Taylor's theorem and is $O(\tau^{3})$ under the same bounds
(cf.\ Jahnke--Lubich~\cite{jahnke2000error},
Hansen--Ostermann~\cite{hansen2009exponential}), so for $\tau\in(0,\tau_0]$
the total defect is bounded by $C_{\mathrm L}\tau^{2}$, as claimed.
\end{proof}

For the operators in \eqref{eq:rho-pde}, a direct computation of the
highest-order ($D\Delta$) contribution to \eqref{eq:bracket} gives
\[
[\mathcal A,\mathcal B](\rho)\Big|_{\text{diff.--react.}}
=D\Big(\Delta\big(\tilde r(\rho)\rho\big)
-\big(\tilde r'(\rho)\rho+\tilde r(\rho)\big)\Delta\rho\Big)
=D\big(\tilde r''(\rho)\rho+2\tilde r'(\rho)\big)\,|\nabla\rho|^{2},
\]
with the advection--reaction cross terms contributing analogous lower-order
(in $D$) factors of $\nabla\rho$. Hence the local splitting error scales with
$|\nabla\rho|^{2}$ and is largest precisely where the solution develops steep
gradients---the diffuse interface of the Allen--Cahn model and the
aggregation core of the Keller--Segel model.
\begin{lemma}\label{lem:stability}
Under Assumption~\ref{ass:regularity}, there exist $\tau_0>0$ and a constant
$L=L\big(D,\|\tilde a\|_{C^{1}_b},\|\tilde r\|_{C^{1}_b},R\big)$ such that the
Lie--Trotter step $\mathcal S_{\tau}=\Phi^{\tau}_{\mathcal B}\circ\Phi^{\tau}_{\mathcal A}$
satisfies
\[
\|\mathcal S_{\tau}\phi-\mathcal S_{\tau}\psi\|_{L^{2}}\le e^{L\tau}\|\phi-\psi\|_{L^{2}},
\qquad \tau\in(0,\tau_0],
\]
for all $\phi,\psi$ with $\|\phi\|_{H^{p}},\|\psi\|_{H^{p}}\le R$.
\end{lemma}

\begin{proof}
We treat the two sub-flows separately and compose the resulting Lipschitz
bounds.
Let $\rho_1(t)=\Phi^{t}_{\mathcal A}\phi$ and $\rho_2(t)=\Phi^{t}_{\mathcal A}\psi$ solve
$\partial_t\rho=\nabla\!\cdot(\tilde a(\rho)\rho)+D\Delta\rho$, and set
$w=\rho_1-\rho_2$. Then $w(0)=\phi-\psi$ and
\[
\partial_t w
=\nabla\!\cdot\!\big(\tilde a(\rho_1)\rho_1-\tilde a(\rho_2)\rho_2\big)
+D\Delta w .
\]
Testing against $w$ in $L^{2}(\mathbb T^{d})$ and integrating by parts,
\[
\frac12\frac{d}{dt}\|w\|_{L^{2}}^{2}
=-\big\langle \nabla w,\ \tilde a(\rho_1)\rho_1-\tilde a(\rho_2)\rho_2\big\rangle
-D\|\nabla w\|_{L^{2}}^{2}.
\]
Decompose the flux difference as
\[
\tilde a(\rho_1)\rho_1-\tilde a(\rho_2)\rho_2
=\tilde a(\rho_1)\,w+\big(\tilde a(\rho_1)-\tilde a(\rho_2)\big)\rho_2 .
\]
Under Assumption~\ref{ass:regularity}, $\|\tilde a\|_{L^\infty}\le C_{b,\mathrm{lip}}$,
$\tilde a$ is Lipschitz with constant $\|\tilde a\|_{C^{1}_b}$, and
 $\|\rho_2\|_{L^\infty}\le C_R$ (from $\|\rho_2\|_{H^{p}}\le R$ and
$H^{p}\hookrightarrow L^\infty$). Hence by Cauchy--Schwarz,
\[
\big|\big\langle \nabla w,\ \tilde a(\rho_1)\rho_1-\tilde a(\rho_2)\rho_2\big\rangle\big|
\le \big(C_{b,\mathrm{lip}}+\|\tilde a\|_{C^{1}_b}C_R\big)\,
\|\nabla w\|_{L^{2}}\,\|w\|_{L^{2}}
=: C_a\,\|\nabla w\|_{L^{2}}\,\|w\|_{L^{2}} .
\]
By Young's inequality, $C_a\|\nabla w\|_{L^2}\|w\|_{L^2}
\le D\|\nabla w\|_{L^2}^2+\tfrac{C_a^2}{4D}\|w\|_{L^2}^2$, so the gradient term is
absorbed by the diffusion term:
\[
\frac12\frac{d}{dt}\|w\|_{L^{2}}^{2}
\le -D\|\nabla w\|_{L^{2}}^{2}+D\|\nabla w\|_{L^{2}}^{2}
+\frac{C_a^{2}}{4D}\|w\|_{L^{2}}^{2}
=\frac{C_a^{2}}{4D}\|w\|_{L^{2}}^{2}.
\]
Thus $\frac{d}{dt}\|w\|_{L^2}^2\le 2L_{\mathcal A}\|w\|_{L^2}^2$ with
$L_{\mathcal A}:=C_a^{2}/(4D)$, and Gr\"onwall's inequality gives
\[
\|w(\tau)\|_{L^{2}}^{2}\le e^{2L_{\mathcal A}\tau}\|w(0)\|_{L^{2}}^{2},
\quad\text{i.e.}\quad
\big\|\Phi^{\tau}_{\mathcal A}\phi-\Phi^{\tau}_{\mathcal A}\psi\big\|_{L^{2}}
\le e^{L_{\mathcal A}\tau}\|\phi-\psi\|_{L^{2}},
\]
with $L_{\mathcal A}=L_{\mathcal A}(D,\|\tilde a\|_{C^{1}_b},C_R)$.

Let now $\rho_1,\rho_2$ solve $\partial_t\rho=\tilde r(\rho)\rho$ and again
$w=\rho_1-\rho_2$. Since $\mathcal B$ involves no spatial derivatives, testing
against $w$ gives
\[
\frac12\frac{d}{dt}\|w\|_{L^{2}}^{2}
=\big\langle w,\ \tilde r(\rho_1)\rho_1-\tilde r(\rho_2)\rho_2\big\rangle .
\]
Writing $\tilde r(\rho_1)\rho_1-\tilde r(\rho_2)\rho_2
=\tilde r(\rho_1)w+(\tilde r(\rho_1)-\tilde r(\rho_2))\rho_2$ and using
$\|\tilde r\|_{L^\infty}\le C_{b,\mathrm{lip}}$, the Lipschitz bound
$\|\tilde r\|_{C^{1}_b}$, and $\|\rho_2\|_{L^\infty}\le C_R$,
\[
\frac12\frac{d}{dt}\|w\|_{L^{2}}^{2}
\le \big(C_{b,\mathrm{lip}}+\|\tilde r\|_{C^{1}_b}C_R\big)\|w\|_{L^{2}}^{2}
=: L_{\mathcal B}\,\|w\|_{L^{2}}^{2}.
\]
Gr\"onwall then yields
$\big\|\Phi^{\tau}_{\mathcal B}\phi-\Phi^{\tau}_{\mathcal B}\psi\big\|_{L^{2}}
\le e^{L_{\mathcal B}\tau}\|\phi-\psi\|_{L^{2}}$ with
$L_{\mathcal B}=L_{\mathcal B}(\|\tilde r\|_{C^{1}_b},C_R)$.
Since $\mathcal S_{\tau}=\Phi^{\tau}_{\mathcal B}\circ\Phi^{\tau}_{\mathcal A}$, the two
estimates compose:
\[
\|\mathcal S_{\tau}\phi-\mathcal S_{\tau}\psi\|_{L^{2}}
\le e^{L_{\mathcal B}\tau}\,
\big\|\Phi^{\tau}_{\mathcal A}\phi-\Phi^{\tau}_{\mathcal A}\psi\big\|_{L^{2}}
\le e^{(L_{\mathcal A}+L_{\mathcal B})\tau}\|\phi-\psi\|_{L^{2}},
\]
which is the claim with $L=L_{\mathcal A}+L_{\mathcal B}$.
\end{proof}


\begin{proposition}[Global splitting error]\label{prop:global-split}
Under Assumption~\ref{ass:regularity}, suppose the time-discrete split solution
$\rho^{n+1}=\mathcal S_{\tau}\rho^{n}$ remains uniformly bounded,
$\sup_{0\le n\le T/\tau}\|\rho^{n}\|_{H^{p}}\le R$. Then, for $\tau\in(0,\tau_0]$,
\[
\max_{0\le n\le T/\tau}\big\|\rho(t_{n})-\rho^{n}\big\|_{L^{2}}
\le \frac{C_{\mathrm L}}{L}\big(e^{LT}-1\big)\,\tau \;=\;O(\tau),
\]
where $C_{\mathrm L}$ is the local-error constant from
Lemma~\ref{lem:local-split} and $L$ is the stability constant from
Lemma~\ref{lem:stability}.
\end{proposition}
\begin{proof}
Set $e_{n}=\rho(t_{n})-\rho^{n}$, so that $e_{0}=0$. By
Assumption~\ref{ass:regularity} the exact solution satisfies
$\rho(t_n)\in\mathcal B_{R}$ for all $t_n\le T$, and by the a priori assumption
$\rho^{n}\in\mathcal B_{R}$ for all $0\le n\le T/\tau$; hence
Lemmas~\ref{lem:local-split}--\ref{lem:stability} apply at every step.

Splitting the exact flow over one step and inserting the local-error
decomposition of Lemma~\ref{lem:local-split},
\[
\rho(t_{n+1})=\Phi^{\tau}\rho(t_{n})
=\mathcal S_{\tau}\rho(t_{n})+d_{n},
\qquad \|d_{n}\|_{L^{2}}\le C_{\mathrm L}\,\tau^{2},
\]
while the scheme advances as $\rho^{n+1}=\mathcal S_{\tau}\rho^{n}$. Subtracting
and applying the $L^{2}$-stability of $\mathcal S_{\tau}$
(Lemma~\ref{lem:stability}) to the pair $\rho(t_n),\rho^{n}\in\mathcal B_{R}$,
\begin{equation}
\|e_{n+1}\|_{L^{2}}
=\big\|\mathcal S_{\tau}\rho(t_n)-\mathcal S_{\tau}\rho^{n}+d_n\big\|_{L^{2}}
\le \|\mathcal S_{\tau}\rho(t_n)-\mathcal S_{\tau}\rho^{n}\|_{L^{2}}+\|d_n\|_{L^{2}}
\le e^{L\tau}\|e_{n}\|_{L^{2}}+C_{\mathrm L}\,\tau^{2}.
\label{eq:onestep}
\end{equation}
Iterating \eqref{eq:onestep} from $e_{0}=0$, each local
defect $d_{k}$ is propagated by $n-1-k$ further stable steps, contributing a
factor $e^{L(n-1-k)\tau}$:
\[
\|e_{n}\|_{L^{2}}
\le \sum_{k=0}^{n-1} e^{L(n-1-k)\tau}\,C_{\mathrm L}\tau^{2}
= C_{\mathrm L}\tau^{2}\sum_{j=0}^{n-1} e^{Lj\tau}
= C_{\mathrm L}\tau^{2}\,\frac{e^{Ln\tau}-1}{e^{L\tau}-1},
\]
where the second equality reindexes $j=n-1-k$ and the third sums the finite
geometric series. Using $n\tau\le T$ and the elementary bound
$e^{L\tau}-1\ge L\tau$ for $\tau>0$,
\[
\|e_{n}\|_{L^{2}}
\le C_{\mathrm L}\tau^{2}\,\frac{e^{LT}-1}{L\tau}
= \frac{C_{\mathrm L}}{L}\big(e^{LT}-1\big)\,\tau,
\qquad 0\le n\le T/\tau,
\]
which is the claimed $O(\tau)$ bound.
\end{proof}


\end{document}
